\subsection{The reciprocal-discriminant sum at the level prime}
\label{app:pair-count}
Write $h=\varrho+1$, $H_0=b^{h^3 f}$ and $\tau=b^{kf/4}$.
We separate a lattice estimate in \S\ref{B:pairs} from the multiplicity
with which a near term can lie over a lattice. Let $\Lambda$ be any lattice
in the genus of the standard definite Hermitian lattice of rank $k$ over
$\mathcal O_E$. At finite places it is self-dual. For a rank-$s$ lattice
$A$ in its rational span, write $A^\#$ for its Hermitian dual there and
$d(A)=[A^\#:A]^{1/2}$ when $A$ is integral. These are cardinality
indices. This is the positive square root of the finite index of the
Gram determinant, the discriminant convention for the pair count. In
particular $d(B)=i\,d(A)$ for $B\subset A$ of index $i$.

\begin{lemma}\label{lem:pair-row-count}
Let $s+j=k$. For each row coset $y$ at the level prime let $t_y$ be its
integral row on $\Lambda$ modulo the ideal $\mathfrak n$ of norm $N^f$.
Suppose the $t_y$ are distinct. For each $y$ take all rank-$s$ sublattices
$B\subset\Lambda$ with $t_y(B)=0$ and $d(B)\le H_0$. At the scales
of the filling their combined weight satisfies
\[
 \sum_y\sum_B d(B)^{-k}
 \le (N^f)^j H_0^{2\varrho+6}
 (200fk\delta_E)^{f(k+1)j}\,2^k .
\]
Here $\delta_E$ is the root discriminant and $(\varrho+1)2f\le2^{2\varrho}$
is the field-degree bound for the filling.
\end{lemma}
\begin{proof}
Put $V=EB$, $A=\Lambda\cap V$, $d=d(A)$, and $i=[A:B]$.
The number of all rows modulo $\mathfrak n$ annihilating $B$ is at most
$i(N^f)^j$. Indeed, locally at $\mathfrak p^m\Vert\mathfrak n$ extend a
basis of the saturated summand $A_{\mathfrak p}$ to an ambient basis,
and diagonalize $B_{\mathfrak p}$ inside it with elementary-divisor
exponents $a_1,\ldots,a_s$. For $q=\operatorname{N}\mathfrak p$
there are $q^{mj+\sum_a\min(a_a,m)}$ choices of row. Multiplication over
the primes gives the assertion, also for a lattice in a moved basis.
Here the rows for the actual row cosets are the distinguished rows of
$y^{-1}$ in coordinates on the unimodular lattice. Congruent inverse
rows put $y^{-1}y'$ in the row stabilizer and give the same coset.

Put $\mu=2\varrho+4$. Since $di\le H_0$ and $i,d\ge1$,
\[
 i(di)^{-k}\le H_0^{\mu+2} d^{-(k+1)} i^{-(s+\mu)}.
\]
For fixed $V$, the sum of $i^{-(s+\mu)}$ over its finite-index sublattices
is at most $2^k$. At a finite place the series for sublattices in a
free module of rank $s$ is $\prod_{a=0}^{s-1}(1-q^{a-z})^{-1}$:
triangular bases over the valuation ring give one diagonal exponent
at each pivot and residues above the pivots. The Euler product at
$z=s+\mu$ is bounded by $\zeta_E(\mu)^s$.
Indeed $\zeta_E(\mu)\le\zeta(\mu)^{2f}$, and
$\log\zeta(\mu)\le 2^{-\mu}(1+2/(\mu-1))$ by integral comparison.
The displayed field-degree bound makes $2f\log\zeta(\mu)<\log 2$.

For clarity, the remaining sum over subspaces satisfies
\[
 \sum_{\operatorname{codim} V=j} d(\Lambda\cap V)^{-(k+1)}
 \le (200fk\delta_E)^{f(k+1)j}.
\]
The case $j=0$ follows from self-duality. For $j>0$ let
$A'=\Lambda\cap V^\perp$. Saturation and self-duality identify
$A^\#/A$ and $(A')^\#/A'$ in cardinality: restriction of the pairing
maps $\Lambda$ onto $A^\#$ locally, since $A$ is a direct summand
locally, and its kernel is $A'$. Thus both quotients have the cardinality
of $\Lambda/(A+A')$. In the Euclidean norm summing the squared lengths
over the $f$ complex places, the covolume of $A'$ in its real span is
$(2^{-f}\sqrt{|D_E|})^j d(A)$. To see the normalization, take a
pseudo-basis of $A'$. The determinant of its complex Gram matrix at
each chosen embedding enters the real covolume to the first power.
The coefficient ideals enter by their absolute norms. Their product
with the Gram factors is $[(A')^\#:A']^{1/2}$, by taking determinants
of the Hermitian pairing locally. The remaining factors are the
Minkowski covolumes of $j$ copies of $\mathcal O_E$.

Let $m=2fj$ and let $\lambda_a$ be the Euclidean successive minima of
$A'$. Minkowski's second theorem, using the cube of side $2/\sqrt m$
inside the unit ball, gives $\prod_{a=1}^m\lambda_a\le
m^{m/2}\operatorname{covol}(A')$. One can choose $E$-independent vectors
$X_1,\ldots,X_j$ with $\|X_{a+1}\|\le\lambda_{2fa+1}$ for $0\le a<j$:
the preceding $E$-span has real dimension $2fa$. Therefore
\[
 (\prod_a \|X_a\|)^{2f}\le (8fj\delta_E)^{fj}d(A).
\]
The chosen tuple determines $V$ by orthogonality. Every nonzero
$x\in\Lambda$ has $\|x\|\ge1$ since the trace from $E^+$ of its Hermitian
norm is a positive integer. The points are thus $1$-separated. Put
$D=2fk$ and $p=2f(k+1)$. Disjoint balls of radius $1/2$ give
$\#\{0<\|x\|\le R\}\le(1+2R)^D$ for $R\ge1$.
Integration against $pR^{-p-1}\,dR$ bounds $\sum_{x\ne0}\|x\|^{-p}$
by $3^D p/(p-D)=3^D(k+1)\le5^p$
(the last inequality follows from $k+1\le(5/3)^{2k}$).
Sum the product bound over all tuples and use $j\le k$ and
$8\cdot5^2=200$. This proves the subspace sum and, on multiplying
the preceding bounds, the lemma.
\end{proof}

\paragraph{Class and pair multiplicities.}
For a free source level the classes have equal quotient measure.
Suppose its row quotient has size $Y\ge(N^k/b)^f$, and the translated
class lattice is identified by a rational unitary isometry with the
kernel of the corresponding row in $\Lambda$ modulo $\mathfrak n$.
Integrate the uniform row bound of Lemma~\ref{lem:pair-row-count}
over the compact adelic quotient with probability measure. Right
translation preserves the measure, so division by $Y$ bounds the
average over classes. At the fitted scales
\[
 b^f H_0^{2\varrho+6}2^k\le\tau,\qquad
 N^f(200fk\delta_E)^{f(k+1)}\le\tau .
\]
For the first comparison use $k\le f$, $b\ge2$ and
$2+h^3(2\varrho+6)\le k/4$. For the second use $N\le b^{60h}$,
$(200fk\delta_E)^{1000}\le b$ and $60h+(k+1)/1000\le k/4$.
Consequently the sum of $d(B)^{-k}$ over classes and their eligible
lattices is at most $\#\{c\}\tau^{j+1}/N^{kf}$.

Finally suppose each near pair determines a common label and a rank-$(k-j)$
common lattice of discriminant at most $H_0$ contained in both term lattices, with pair discriminant
equal to the discriminant of the common lattice. The single-term estimate
needed is that, at fixed first class, label and such lattice, there are at
most $b^{2000h^6(j+1)f}$ possible terms, counting the second class as
part of the term. Thus the number of ordered pairs in a fibre is at most
the square of this bound. For $j\le h^3$ and $k\ge512000 h^9$ that square
is at most $\tau$. Summing over labels and multiplying by the preceding
class sum gives $C_p\tau^{j+2}$, where
$C_p=\#\{\text{labels}\}\#\{c\}/N^{kf}$.
This proves the reciprocal-discriminant reduction in \S\ref{B:pairs}
given the indicated fitting data and single-term multiplicity. The use of this
bounded-discriminant count for all the near pairs requires the geometric
bound $d\le H_0$ in addition to $j\le h^3$; the source derives these together
in the \lineSrc{New/Spanning/RemainingTerms/Sum}{90}{gauge comparison}. The local
cell-rank and gauge estimates that bound the weights of individual pairs
are separate inputs. The
\lineSrc{New/Spanning/RemainingTerms/LatticesOfCorank}{373}{row count},
\lineSrc{New/Spanning/RemainingTerms/SiegelLemma}{630}{subspace sum} and
\lineSrc{New/Spanning/RemainingTerms/Count}{187}{pair estimate}
specify the same normalizations.
